% SPDX-License-Identifier: CC-BY-SA-4.0
% Author: Matthieu Perrin

\begingroup

\begin{frame}[fragile]{La méthode \lstinline{enqueue}}

  \vspace{-20mm}
  \begin{onlyenv}<1>
    \begin{lstlisting}[gobble=6]
      public void enqueue(T value) {





        (*\Alert{last.next}*) = new Node(value);
        (*\Alert{last}*) = last.next;




      }
    \end{lstlisting}
  \end{onlyenv}
  \begin{onlyenv}<2>
    \begin{lstlisting}[gobble=6]
      public void enqueue(T value) {





        last. (*\Alert{get()}*) .next. (*\Structure{set}*) (new Node(value));
        last. (*\Structure{set}*) (last. (*\Alert{get()}*) .next. (*\Example{get()}*));




      }
    \end{lstlisting}
  \end{onlyenv}
  \begin{onlyenv}<3>
    \begin{lstlisting}[gobble=6]
      public void enqueue(T value) {
        var node = new Node(value);
        
        var l = last. (*\Alert{get();}*)


        l.next. (*\Structure{set(node);}*)
        last. (*\Structure{set(node);}*)




      }
    \end{lstlisting}
  \end{onlyenv}
  \begin{onlyenv}<4>
    \begin{lstlisting}[gobble=6]
      public void enqueue(T value) {
        var node = new Node(value);
        
        var l = last.get();


        l.next. (*\Structure{compareAndSet(null, node);}*)
        last. (*\Structure{compareAndSet(l, node);}*)




      }
    \end{lstlisting}
  \end{onlyenv}
  \begin{onlyenv}<5>
    \begin{lstlisting}[gobble=6]
      public void enqueue(T value) {
        var node = new Node(value);
        (*\Structure{while (true)}*) {
          var l = last.get();


          (*\Structure{if (l.next.compareAndSet(null, node))}*) {
            last.compareAndSet(l, node);
            (*\Structure{return;}*)
          }
          
        }
      }
    \end{lstlisting}
  \end{onlyenv}
  \begin{onlyenv}<6>
    \begin{lstlisting}[gobble=6]
      public void enqueue(T value) {
        var node = new Node(value);
        while (true) {
          var l = last.get();
          (*\Structure{var n = l.next.get();}*)
          (*\Structure{if (n == null)}*) {
            if (l.next.compareAndSet(n, node)){
              last.compareAndSet(l, node);
              return;
            }
          } (*\Structure{else last.compareAndSet(l, n);}*)
        }
      }
    \end{lstlisting}
  \end{onlyenv}

  \onBlock<3->[right=45mm, top]{Conseils}{
    \begin{enumerate} \setcounter{enumi}{3}
    \item<3-> Une seule opération atomique par ligne
    \item<4-> Protéger les écritures par des \lstinline{compareAndSet}
    \item<5-> Retenter si échec
      \begin{itemize}
      \item lock-freedom
      \end{itemize}
    \end{enumerate}
  }

  \onAlertBlock<3-5>[y=-10mm, anchor=north]{Contre-exemple}{
    \begin{tikzpicture}[anchor=mid, x=15mm]
      \draw<3-4>[process] (0,1) node[left]{$T_1$} to (7,1);
      \draw<3-5>[process] (0,0) node[left]{$T_2$} to (7,0);
      \draw<5>  [crashed] (0,1) node[left]{$T_1$} to (2.6,1);
      
      \node<3>[alert, operation] (31) at (1 ,1) {$get()$};
      \node<3>[alert, operation] (32) at (2 ,0) {$get()$};
      \node<3>[alert, operation] (33) at (3 ,1) {$set()$};
      \node<3>[alert, operation] (34) at (4 ,0) {$set()$};
      \node<3>[alert, operation] (35) at (5 ,0) {$set()$};
      \node<3>[alert, operation] (36) at (6 ,1) {$set()$};
      \node<4>[alert, operation] (41) at (1 ,1) {$get()$};
      \node<4>[alert, operation] (42) at (2 ,0) {$get()$};
      \node<4>[alert, operation] (43) at (3 ,1) {$cas()$};
      \node<4>[alert, operation] (44) at (4 ,0) {$cas()$};
      \node<4>[alert, operation] (45) at (5 ,0) {$cas()$};
      \node<4>[alert, operation] (46) at (6 ,1) {$cas()$};
      \node<5>[alert, operation] (51) at (1 ,1) {$get()$};
      \node<5>[alert, operation] (52) at (2 ,1) {$cas()$};
      \node<5>[alert, operation] (53) at (3 ,0) {$get()$};
      \node<5>[alert, operation] (54) at (4 ,0) {$cas()$};
      \node<5>[alert, operation] (55) at (5 ,0) {$get()$};
      \node<5>[alert, operation] (56) at (6 ,0) {$cas()$};

      \begin{scope}[background]
        \node<3>[structure, operation, fit=(31)(36)] {};
        \node<3>[structure, operation, fit=(32)(35)] {};
        \node<4>[structure, operation, fit=(41)(46)] {};
        \node<4>[structure, operation, fit=(42)(45)] {};
        \node<5>[structure, operation, fit=(51)(52)] {};
        \node<5>[structure, operation, fit=(53)(56)] {};
      \end{scope}
      
    \end{tikzpicture}
  }
  
\end{frame}

\endgroup
\endinput
